
\documentclass{article} % For LaTeX2e
\usepackage{iclr2026_conference,times}

% Optional math commands from https://github.com/goodfeli/dlbook_notation.
\input{math_commands.tex}

\usepackage{hyperref}
\usepackage{url}
\usepackage{booktabs}
\usepackage{multirow}

% Notation macros for the GAPSBI data contract (defined defensively to avoid
% clashing with math_commands.tex).
\providecommand{\xfull}{\mathbf{x}_{\mathrm{full}}}
\providecommand{\xobs}{\mathbf{x}_{\mathrm{obs}}}
\providecommand{\xtilde}{\widetilde{\mathbf{x}}}
\providecommand{\xhat}{\widehat{\mathbf{x}}}
\providecommand{\mvec}{\mathbf{m}}
\providecommand{\thetab}{\boldsymbol{\theta}}
\providecommand{\phib}{\boldsymbol{\phi}}


\title{GAPSBI: A Benchmark for Simulation-Based Inference under Missing Data}

% Authors must not appear in the submitted version. They should be hidden
% as long as the \iclrfinalcopy macro remains commented out below.
% Non-anonymous submissions will be rejected without review.

\author{Antiquus S.~Hippocampus, Natalia Cerebro \& Amelie P. Amygdale \thanks{ Use footnote for providing further information
about author (webpage, alternative address)---\emph{not} for acknowledging
funding agencies.  Funding acknowledgements go at the end of the paper.} \\
Department of Computer Science\\
Cranberry-Lemon University\\
Pittsburgh, PA 15213, USA \\
\texttt{\{hippo,brain,jen\}@cs.cranberry-lemon.edu} \\
\And
Ji Q. Ren \& Yevgeny LeNet \\
Department of Computational Neuroscience \\
University of the Witwatersrand \\
Joburg, South Africa \\
\texttt{\{robot,net\}@wits.ac.za} \\
\AND
Coauthor \\
Affiliation \\
Address \\
\texttt{email}
}

% The \author macro works with any number of authors. There are two commands
% used to separate the names and addresses of multiple authors: \And and \AND.
%
% Using \And between authors leaves it to \LaTeX{} to determine where to break
% the lines. Using \AND forces a linebreak at that point. So, if \LaTeX{}
% puts 3 of 4 authors names on the first line, and the last on the second
% line, try using \AND instead of \And before the third author name.

\newcommand{\fix}{\marginpar{FIX}}
\newcommand{\new}{\marginpar{NEW}}

%\iclrfinalcopy % Uncomment for camera-ready version, but NOT for submission.
\begin{document}


\maketitle

\begin{abstract}
Simulation-based inference (SBI) increasingly targets scientific settings where
observations are incomplete, yet missing data is rarely treated as a controlled,
systematically benchmarked axis of evaluation. We introduce GAPSBI, a benchmark
and experimentation framework for simulation-based inference under missing data.
GAPSBI couples a small set of synthetic simulators with explicit missingness
mechanisms spanning missing-completely-at-random (MCAR), missing-at-random (MAR),
and missing-not-at-random (MNAR) processes, reproducible HDF5 datasets with
fixed train/validation/test splits and seeds, a family of neural posterior
estimation (NPE) baselines that share a common posterior-estimation backbone, and
a standardized evaluation protocol. The protocol combines posterior calibration
diagnostics (simulation-based calibration and TARP coverage) with posterior shift
metrics that compare each missing-data posterior against fixed high-quality
reference posteriors, together with full-data seed-to-seed variability baselines that
quantify how much observed shift is attributable to training stochasticity.
GAPSBI also defines distinct simulation-budget regimes so that method behavior
can be studied as a function of the amount of training simulation. This draft
documents the benchmark design, tasks, missingness mechanisms, simulation
budgets, baseline methods, and evaluation protocol; empirical results are
reported in companion sections.
\end{abstract}

\section{Introduction}
\label{sec:introduction}

Simulation-based inference (SBI) enables Bayesian parameter estimation for models
defined only through a stochastic simulator, by learning surrogate posteriors
from simulated parameter-observation pairs. Neural posterior estimation (NPE),
in particular, trains a conditional density estimator to approximate the
posterior over simulator parameters given an observation. In many scientific
applications, however, the observation is incomplete: sensors drop out, parts of
a time series are unrecorded, or measurements are censored. Practitioners
therefore must adapt SBI methods to partially observed data, but the field lacks
a controlled, reproducible benchmark that isolates the effect of missingness on
posterior quality.

This motivates a central question for GAPSBI: which missing-data representation
works best for a fixed Neural Posterior Estimation (NPE) backbone? Existing
approaches to missing data often change several ingredients at once: the
representation of partial observations, the inference architecture, the
optimization objective, and sometimes the inference procedure itself. Such
comparisons are scientifically useful but difficult to interpret when the goal is
to understand the missing-data representation. GAPSBI instead fixes the
downstream posterior estimator to NPE with a neural spline flow (NSF) density
estimator wherever possible, so that differences between baselines are primarily
attributable to how $\xobs$ and $\mvec$ are represented before posterior
estimation.

GAPSBI addresses this gap. It is a benchmark and experimentation framework for
simulation-based inference under missing data, built around a deliberately simple
and explicit data contract: each dataset stores complete simulator outputs
$\xfull$, a binary mask $\mvec$ where $1$ denotes observed and $0$ denotes
missing, and the observed data $\xobs = \xfull \odot \mvec$, together with the
underlying parameters $\thetab$. Datasets are saved in a shared HDF5 layout with
fixed train, validation, and test splits and explicit random seeds, so that
experiments are reproducible and comparable across methods. The formal setup is
given in Section~\ref{sec:benchmark}.

The benchmark spans several synthetic simulators and three families of missingness
mechanisms (MCAR, MAR, and MNAR), providing a structured space in which to study
how different missingness processes affect inference. On top of the datasets,
GAPSBI provides a set of NPE baselines for handling missingness, ranging from a
full-data no-missingness model to imputation, mask augmentation, masked embeddings,
learned imputation, and a RISE-style probabilistic imputation baseline. It
further provides a standardized evaluation protocol that goes
beyond visual inspection of posterior samples: it measures posterior calibration
with simulation-based calibration (SBC) and TARP coverage, and quantifies how
much each method posterior departs from a fixed high-quality reference posterior using
distributional and moment-based posterior shift metrics. Because NPE training is
stochastic, GAPSBI also computes full-data seed-to-seed variability baselines,
so that observed posterior shifts can be interpreted relative to intrinsic
training variability rather than attributed naively to missingness.

\paragraph{Contributions.} This paper makes the following contributions, all of
which are reflected in the accompanying open-source repository:
\begin{itemize}
  \item A reproducible benchmark framework for SBI under missing data, with a
  simple HDF5 dataset contract, explicit seeds, and standardized train,
  validation, and test splits.
  \item A collection of synthetic simulators together with explicit MCAR, MAR,
  and MNAR missingness mechanisms that can be combined at controlled missingness
  fractions.
  \item A family of NPE baselines for partially observed data: full-data NPE,
  zero- and mean-imputation NPE, mask-augmentation NPE, masked-pooling and
  masked-attention embedding NPE, learned-imputation NPE, and NPE with
  RISE-style probabilistic imputation.
  \item A controlled experimental design that isolates missing-data
  representations while holding the posterior-estimation backbone fixed, enabling
  direct comparisons of how partial observations are encoded for inference.
  \item A standardized evaluation protocol combining calibration diagnostics
  (SBC, TARP), posterior shift metrics relative to high-quality reference
  posteriors, and full-data variability baselines.
  \item A simulation-budget design that defines distinct high-, mid-, and
  low-budget regimes over the same datasets, enabling study of method behavior
  as a function of training simulation budget.
\end{itemize}

\section{Related Work}
\label{sec:related-work}

\textit{Placeholder.} This section will situate GAPSBI within prior work. Planned
coverage (to be completed with appropriate citations):
\begin{itemize}
  \item Simulation-based inference and neural posterior estimation methods.
  Transformer-based SBI methods, including Simformer-style architectures, pursue
  a complementary design philosophy: they modify the inference architecture
  itself to process structured or partially observed inputs. GAPSBI instead
  fixes the posterior estimator and compares missing-data representations under a
  common NPE backbone, separating representation effects from architectural
  effects.
  \item Existing SBI benchmarks and standardized evaluation efforts.
  \item Calibration diagnostics for posterior approximations, including
  simulation-based calibration and coverage-based tests.
  \item Methods for inference and learning under missing data, including
  imputation, masking, and missingness-aware modeling.
\end{itemize}

\section{The GAPSBI Benchmark}
\label{sec:benchmark}

GAPSBI is organized around reproducible datasets, a fixed set of simulators,
explicit missingness mechanisms, and clearly separated simulation-budget regimes.
This section describes the design goals, the benchmark tasks, the missingness
mechanisms, and the simulation budgets.

\paragraph{Problem setup and notation.} Each benchmark instance is defined by a
prior $p(\thetab)$ over parameters $\thetab \in \mathbb{R}^{d_\theta}$ and a
stochastic simulator that induces a likelihood $p(\xfull \mid \thetab)$ over
complete observations $\xfull \in \mathbb{R}^{d}$. A complete example is drawn as
\begin{equation}
  \thetab \sim p(\thetab), \qquad \xfull \sim p(\xfull \mid \thetab).
  \label{eq:generative}
\end{equation}
Missingness is then introduced through a binary mask
$\mvec \in \{0,1\}^{d}$ sampled from a mechanism-dependent distribution, and the
observed data is obtained by elementwise multiplication,
\begin{equation}
  \xobs = \xfull \odot \mvec, \qquad
  m_i =
  \begin{cases}
    1 & \text{if entry $i$ is observed},\\
    0 & \text{if entry $i$ is missing},
  \end{cases}
  \label{eq:contract}
\end{equation}
where $\odot$ denotes elementwise multiplication. Each dataset stores the tuple
$(\thetab, \xfull, \xobs, \mvec)$ for every example, partitioned into train,
validation, and test splits.

\subsection{Design Goals}
\label{sec:design-goals}

GAPSBI is guided by a small number of explicit design goals.

\paragraph{Reproducibility and an explicit data contract.} Every dataset follows
the same contract of Equations~\eqref{eq:generative}--\eqref{eq:contract}:
generate complete simulator outputs $\xfull$, generate a binary mask $\mvec$ with
the convention $1 = \text{observed}$ and $0 = \text{missing}$, store observed data
as $\xobs = \xfull \odot \mvec$, and save train, validation, and test splits in a
shared HDF5 format. Randomness is kept explicit through fixed seeds. Datasets use
a grouped HDF5 layout with \texttt{/train}, \texttt{/val}, and \texttt{/test}
groups, each containing \texttt{theta}, \texttt{x\_full}, \texttt{x\_obs}, and
\texttt{mask} arrays, with simulator and mask metadata stored as file-level
attributes (nested metadata serialized as JSON). This makes datasets
self-describing and directly comparable across methods.

\paragraph{Simplicity and low engineering burden.} The benchmark deliberately
favors small, cheap, and easily reproducible problems and relies on standard NPE
tooling. This keeps the core baselines easy to run and reproduce, including in a
purely local setting, and lowers the barrier to extending the benchmark.

\paragraph{Systematic missingness coverage.} A central goal is to compare MCAR,
MAR, and MNAR missingness carefully and at controlled missingness fractions,
rather than studying a single ad hoc missingness pattern. This is reflected in
the explicit mechanism implementations described in
Section~\ref{sec:missingness}.

\paragraph{Controlled architectural comparisons.} GAPSBI is designed to study
representation choices rather than conflate them with changes in the inference
architecture. Whenever possible, benchmark methods share the same downstream
posterior estimator---an NSF-based NPE model---and differ only in the
conditioning variable $\mathbf{c}$ passed to that estimator. This makes the
comparison focus on how partially observed data are encoded, completed, or
embedded before inference.

\paragraph{Standardized, calibration-aware evaluation.} GAPSBI is built so that
methods are evaluated with a shared protocol that emphasizes calibration (SBC and
TARP) and quantitative posterior comparison, rather than relying on visual
inspection of posterior samples. The full protocol is described in
Section~\ref{sec:evaluation}.

\paragraph{Modularity and extensibility.} The benchmark is organized into
modular task and method directories, so that tasks, representations,
architectures, and analysis can be extended incrementally without disrupting
existing baselines.

\paragraph{Persisted experiment artifacts.} Each run writes not only posterior
samples and diagnostics, but also the trained per-seed NPE checkpoint. The
checkpoint contains the fitted conditional density estimator, preprocessing maps
for $\thetab$ and $\xobs$, run metadata, and any method-specific components such
as an embedding network or imputer. This makes it possible to reconstruct
$q_{\phib}(\thetab \mid \cdot)$ for additional posterior sampling without
retraining.

\subsection{Benchmark Tasks}
\label{sec:tasks}

GAPSBI implements several synthetic simulators. Each simulator defines a parameter
vector $\thetab \in \mathbb{R}^{d_\theta}$, an observation
$\xfull \in \mathbb{R}^{d}$, and a prior over parameters.

\begin{itemize}
  \item \textbf{OUP} (\texttt{oup}): an Ornstein--Uhlenbeck process with
  $\thetab = (\theta_1, \log \theta_2)$, $d_\theta = 2$, and $d=25$. Prior:
  $\theta_1 \sim \mathcal{U}(0,2)$ and
  $\log \theta_2 \sim \mathcal{U}(-2,3)$.
  \item \textbf{GLU} (\texttt{glu}): the Gaussian Linear Uniform benchmark with
  $d_\theta = d = 10$. Prior:
  $\theta_i \sim \mathcal{U}(-1,1)$ for each coordinate; the simulator draws
  $\xfull = \thetab + \boldsymbol{\epsilon}$ with Gaussian observation noise.
  \item \textbf{GLM} (\texttt{glm}): a Bernoulli generalized linear model with
  fixed stimulus features, $d_\theta=10$, and a default duration of $100$ time
  bins. Prior: an SBIBM-style Gaussian prior, with intercept
  $\theta_0 \sim \mathcal{N}(0,2)$ and stimulus-filter weights
  $\thetab_{1:9} \sim \mathcal{N}(\mathbf{0},(F^\top F)^{-1})$. In summary
  mode, $\xfull$ contains a spike count and spike-triggered stimulus summaries;
  in raw mode, $\xfull$ is the binary spike train.
  \item \textbf{Lotka--Volterra} (\texttt{lotka\_volterra}): a two-population
  predator--prey ODE benchmark with
  $\thetab=(\alpha,\beta,\gamma,\delta)$, $d_\theta=4$, and an interleaved
  prey/predator observation vector of length $d=100$ formed from $50$ timestamps.
  Prior: $\log \thetab \sim
  \mathcal{N}((-0.125,-3.0,-0.125,-3.0), 0.5^2 I)$.
  \item \textbf{Spatial SIR} (\texttt{spatial\_sir}): a spatial SIR lattice model
  with $\thetab=(\beta,\gamma)$, $d_\theta=2$, a default $16 \times 16$ lattice,
  and a single final susceptible/infected/recovered snapshot at measurement time
  $0.25$, flattened to $d=768$. Prior:
  $\beta \sim \mathcal{U}(0,1)$ and $\gamma \sim \mathcal{U}(0,1)$.
  \item \textbf{Hodgkin--Huxley} (\texttt{hodgkin\_huxley}): a Hodgkin--Huxley
  neuron model with $\thetab=(g_{\mathrm{Na}}, g_{\mathrm{K}})$,
  $d_\theta=2$, and a downsampled raw voltage trace as the observation
  (approximately $d=601$ by default). Prior:
  $g_{\mathrm{Na}} \sim \mathcal{U}(0.5,80.0)$ and
  $g_{\mathrm{K}} \sim \mathcal{U}(10^{-4},15.0)$.
\end{itemize}

The main benchmark focuses on four lightweight tasks (OUP, GLM, GLU,
and Lotka--Volterra), which are inexpensive enough for repeated missingness,
method, seed, and simulation-budget sweeps; Spatial SIR, Hodgkin--Huxley, and a
legacy Ricker simulator are implemented but are not part of the main benchmark.
Table~\ref{tab:priors}
lists the parameters and priors of these tasks, and
Table~\ref{tab:tasks-overview} summarizes their dimensions and observation types.

\begin{table}[t]
\centering
\caption{Parameters and priors for the main benchmark tasks.
$\mathcal{U}(a,b)$ denotes a uniform distribution on $[a,b]$.}
\label{tab:priors}
\begin{tabular}{@{}lp{0.25\linewidth}p{0.50\linewidth}@{}}
\toprule
Problem & Parameter & Prior \\
\midrule
\multirow{2}{*}{OUP}
  & $\theta_1$       & $\mathcal{U}(0, 2)$ \\
  & $\log \theta_2$  & $\mathcal{U}(-2, 3)$ \\
\midrule
\multirow{2}{*}{GLM}
  & $\theta_0$ & $\mathcal{N}(0,2)$ \\
  & $\thetab_{1:9}$ & $\mathcal{N}(\mathbf{0}, (F^\top F)^{-1})$ \\
\midrule
GLU & $\theta_i,\ i = 1, \dots, 10$ & $\mathcal{U}(-1, 1)$ \\
\midrule
Lotka--Volterra & $(\alpha,\beta,\gamma,\delta)$ &
$\log \alpha,\log \gamma \sim \mathcal{N}(-0.125,0.5^2)$;
$\log \beta,\log \delta \sim \mathcal{N}(-3.0,0.5^2)$ \\
\bottomrule
\end{tabular}
\end{table}

\begin{table}[t]
\centering
\caption{Overview of the main benchmark tasks: parameter dimension
$\dim(\thetab)$, observation dimension $\dim(\xfull)$, and observation type. GLM
is used in its raw spike-train configuration (duration $100$).}
\label{tab:tasks-overview}
\begin{tabular}{@{}lccl@{}}
\toprule
Problem & $\dim(\thetab)$ & $\dim(\xfull)$ & Observation type \\
\midrule
OUP    & $2$  & $25$  & Ornstein--Uhlenbeck time series \\
GLM    & $10$ & $100$ & Binary spike train \\
GLU    & $10$ & $10$  & Gaussian observation vector \\
Lotka--Volterra & $4$ & $100$ & Interleaved prey/predator time series \\
\bottomrule
\end{tabular}
\end{table}

\subsection{Missingness Mechanisms}
\label{sec:missingness}

GAPSBI implements three families of missingness mechanisms, all using the
convention that a mask value of $1$ denotes an observed entry and $0$ a missing
entry, with $\xobs = \xfull \odot \mvec$. The families differ in what the mask
distribution $p(\mvec \mid \xfull, \thetab)$ is allowed to depend on, which is the
standard distinction between MCAR, MAR, and MNAR. We write $p$ for the target
missing fraction.

\paragraph{MCAR.} Missing-completely-at-random masks are independent of both the
observed values and index metadata,
\begin{equation}
  p(\mvec \mid \xfull, \thetab) = p(\mvec).
  \label{eq:mcar}
\end{equation}
\texttt{PointMCARMask} applies independent Bernoulli masking at each entry,
$m_i \sim \mathrm{Bernoulli}(1 - p)$, while \texttt{BlockMCARMask} masks
contiguous blocks for one-dimensional or batched one-dimensional data.

\paragraph{MAR.} \texttt{CoordinateMARMask} implements coordinate- or
time-dependent missingness whose probability depends only on index metadata $i$
along the last axis, not on the values of $\xfull$,
\begin{equation}
  p(\mvec \mid \xfull, \thetab) = p(\mvec \mid i),
  \label{eq:mar}
\end{equation}
which makes it missing-at-random rather than MNAR. It supports
\texttt{increasing}, \texttt{decreasing}, and \texttt{middle} modes; the requested
missing fraction is approximately the realized missing fraction before probability
clipping, and the realized fraction can be lower when the maximum probability
clips values.

\paragraph{MNAR.} \texttt{SelfCensoringMNARMask} implements value-dependent
self-censoring, so the mask distribution depends on $\xfull$ through a per-entry
score $\mathbf{s} = \mathrm{normalize}(g(\xfull)) \in [0,1]^{d}$,
\begin{equation}
  p(\mvec \mid \xfull, \thetab) = p(\mvec \mid \mathbf{s}),
  \qquad \mathbf{s} = \mathrm{normalize}(g(\xfull)),
  \label{eq:mnar}
\end{equation}
where each sample is optionally transformed by $g$, min/max shifted into
$[0,1]$, and higher normalized scores $s_i$ receive higher missingness
probability; constant samples use a score of $0.5$ everywhere to avoid division
instability. The transform $g$ is one of \texttt{identity} ($g(\xfull) = \xfull$),
\texttt{log1p} ($g(\xfull) = \log(1 + \xfull)$, suited to nonnegative spiky or
count data), or \texttt{abs} ($g(\xfull) = |\xfull|$, suited to
signed vector data where magnitude should drive missingness).
For OUP, GLM, and GLU, the main benchmark uses a mean-normalized variant of this
score, which keeps the value-dependent self-censoring behavior while making the
realized missing fractions closer to the requested nominal levels. Because
probabilities are clipped, high-MNAR settings can still realize less missingness
than the nominal fraction; we therefore report realized mask statistics.

For tasks with structured observations such as Spatial SIR, masking is applied at
the spatial cell level: a single cell-level mask is expanded across the
susceptible, infected, and recovered channels before flattening, so all channels
for a cell are jointly observed or missing.

For Lotka--Volterra, the two population channels are packed into a one-dimensional
interleaved vector. Missingness is therefore applied at the timestamp level: a
single MCAR, MAR, or MNAR decision observes or hides both prey and predator at the
same time point. Its MNAR mechanism uses the log total population
$\log(\mathrm{prey}_t+\mathrm{predator}_t)$ as the timestamp-level score.

\subsection{Simulation Budgets}
\label{sec:budgets}

GAPSBI separates the amount of training simulation into explicit budget regimes
that share the same underlying datasets, seeds, and test split, so that method
behavior can be compared across simulation budgets without confounds.

The canonical dataset configuration uses fixed split sizes of $90{,}000$ training,
$10{,}000$ validation, and $1{,}000$ test examples, generated with dataset seed
$123$. Missingness is applied at fractions $0.10$, $0.25$, and $0.50$. The
canonical mechanism instantiations are point MCAR, coordinate MAR with
\texttt{mar-mode = increasing}, and mean-normalized identity self-censoring MNAR
for OUP, GLM, and GLU. Lotka--Volterra uses timestamp-level analogues: block
MCAR, increasing time-index MAR, and log-total-population MNAR.

On top of these datasets, configurations are organized into three
simulation-budget regimes:
\begin{itemize}
  \item \textbf{High simulation budget}: the canonical
  $90$k\,/\,$10$k\,/\,$1$k train/validation/test setup.
  \item \textbf{Mid simulation budget}: a $9$k train / $1$k validation ablation
  that reuses the same HDF5 datasets and the full $1{,}000$-example test split.
  \item \textbf{Low simulation budget}: a $900$ train / $100$ validation
  ablation that also reuses the same datasets and test split.
\end{itemize}

Within each budget regime, every configuration is trained with a fixed list of
five random seeds, allowing seed-level aggregation and variability analysis (see
Section~\ref{sec:evaluation}).

\section{Methods}
\label{sec:methods}

All benchmark methods use the same downstream NPE density estimator, an NSF
conditional flow, and differ only in the conditioning variable $\mathbf{c}$ or in
the representation used to construct it from $\xfull$, $\xobs$, and $\mvec$.
This isolates the scientific question of missing-data representation: when
$q_{\phib}(\thetab \mid \mathbf{c})$ is held fixed, which encoding of partial
observations yields the best calibrated and most faithful posterior? Unless
stated otherwise, observations are transformed using scalers fit on the training
split, and all posterior sampling is performed from the trained conditional
density estimator.

\subsection{Full-Data NPE}
\label{sec:method-full-data}

The full-data baseline trains
\begin{equation}
  q_{\phib}(\thetab \mid \mathbf{c}), \qquad \mathbf{c} = T_x(\xfull),
  \label{eq:full-data-npe}
\end{equation}
where $T_x$ denotes the training-set preprocessing transform. This baseline is
not a missing-data method; it provides a no-missingness NPE baseline and a
seed-to-seed variability reference, but not the final gold-standard posterior
used for fidelity metrics.

\subsection{Imputation Baselines}
\label{sec:method-imputation}

The simplest missing-data baselines first construct a completed observation
$\xtilde$ and then train
\begin{equation}
  q_{\phib}(\thetab \mid T_x(\xtilde)).
  \label{eq:imputation-npe}
\end{equation}
For zero imputation, missing entries are set to $0$ in transformed observation
space. For mean imputation, a missing coordinate $i$ is set to the empirical
training mean of coordinate $i$ computed from observed entries only. These
baselines intentionally discard explicit information about which entries were
missing.

\subsection{Mask Augmentation}
\label{sec:method-mask-augmentation}

Mask augmentation exposes the missingness pattern to the posterior estimator.
After zero-imputing in transformed observation space, the conditioning vector is
\begin{equation}
  \mathbf{c} = [\,T_x(\xtilde),\, \mvec\,],
  \label{eq:mask-augmentation}
\end{equation}
and NPE trains $q_{\phib}(\thetab \mid \mathbf{c})$. This lets the estimator use
both the observed values and the event that particular coordinates are missing,
which is especially relevant when the missingness mechanism carries information.

\subsection{Masked Embedding NPE}
\label{sec:method-masked-embedding}

The masked-embedding baselines use the same pair
$[\,T_x(\xtilde), \mvec\,]$ but pass it through a learned embedding
$e_\psi$ before density estimation:
\begin{equation}
  q_{\phib}(\thetab \mid e_\psi(T_x(\xtilde), \mvec)).
  \label{eq:masked-embedding}
\end{equation}
The masked-pooling variant embeds each coordinate as a token, pools only observed
tokens, and appends a summary of $\mvec$. The masked-attention variant applies a
small self-attention encoder in which missing coordinates are excluded as
attention keys and values while the mask summary remains available to the
estimator.

\subsection{Learned Imputation}
\label{sec:method-learned-imputation}

Learned imputation replaces fixed completion rules with an imputer
$g_\psi(T_x(\xobs), \mvec)$ trained jointly with NPE. The completed observation is
\begin{equation}
  \xhat
  =
  \mvec \odot T_x(\xobs)
  +
  (1 - \mvec) \odot g_\psi(T_x(\xobs), \mvec),
  \label{eq:learned-imputation}
\end{equation}
and the posterior estimator is trained as $q_{\phib}(\thetab \mid \xhat)$. The
training objective combines the NPE negative log likelihood with a reconstruction
loss on missing coordinates, using the stored $\xfull$ available in simulated
training data.

\subsection{NPE RISE-style Probabilistic Imputation}
\label{sec:method-npe-rise}

The NPE-RISE baseline uses a probabilistic imputer that predicts a completion
distribution rather than only a point completion. Given $(T_x(\xobs), \mvec)$,
the imputer outputs coordinate-wise Gaussian parameters
$(\boldsymbol{\mu}_\psi, \boldsymbol{\sigma}_\psi)$ and forms
\begin{equation}
  \xhat
  =
  \mvec \odot T_x(\xobs)
  +
  (1 - \mvec) \odot \boldsymbol{\mu}_\psi(T_x(\xobs), \mvec).
  \label{eq:npe-rise-completion}
\end{equation}
The density estimator is trained on $q_{\phib}(\thetab \mid \xhat)$ together with
auxiliary imputation losses. An optional mask-prediction head is enabled for MNAR
settings, allowing the imputer to model information contained in the missingness
process itself.

\section{Evaluation Protocol}
\label{sec:evaluation}

GAPSBI evaluates methods with a standardized protocol applied to held-out test
observations. For each configuration (task, missingness mechanism, missingness
fraction, method, and simulation budget) and each of the five seeds, a model is
trained and posteriors are sampled on the test split. Each run produces per-seed
artifacts, including the trained model checkpoint
(\texttt{model\_checkpoint.pt}), posterior samples
(\texttt{posterior\_samples.h5}), calibration plots (\texttt{tarp.png},
\texttt{sbc\_rank\_histograms.png}), diagnostic arrays
(\texttt{diagnostics\_arrays.npz}), and a run summary (\texttt{summary.json}).
The checkpoint stores the fitted approximation $q_{\phib}(\thetab \mid \cdot)$,
the preprocessing transformations used to construct its conditioning variables,
and any learned imputation or embedding components needed to reconstruct the
model for additional inference. Seed-level summaries are aggregated into a master
results table for analysis, and stability across seeds and configurations is
summarized separately.

\subsection{Reference Posteriors}
\label{sec:eval-reference-posteriors}

For posterior fidelity metrics, GAPSBI uses a fixed set of ten reference
observations per main benchmark task, stored separately from the training datasets.
Each reference artifact contains the unscaled parameter draw $\thetab^\star$, the
unscaled complete observation $\xfull^\star$, and $10{,}000$ unscaled posterior
samples from a high-quality task-specific reference posterior
$p(\thetab \mid \xfull^\star)$. GLU uses its analytic truncated-Gaussian
posterior; OUP uses a high-resolution deterministic grid posterior; GLM uses
emcee MCMC under the SBIBM-style Gaussian prior; and Lotka--Volterra uses emcee
MCMC in log-parameter space under the lognormal prior. These reference
posteriors are validated with marginal and joint posterior plots, posterior
predictive checks, and, for MCMC references, convergence diagnostics including
split $\widehat{R}$, effective sample size, trace plots, and independent
ensemble checks.

Calibration diagnostics use the held-out test split because SBC and TARP require
many independent parameter-observation pairs. Fidelity and posterior-shift
metrics instead use the ten fixed reference observations, following the standard
benchmarking pattern of evaluating approximate posteriors against a smaller
number of carefully constructed reference posteriors.

For missing-data methods, the ten complete reference observations are converted
to method inputs by applying the same missingness mechanism and fraction as the
corresponding experiment, using deterministic reference masks, and then applying
the method-specific preprocessing or imputation rule. These method posteriors
are compared to the full-observation reference posteriors. Thus, for missing-data
methods, including imputation and mask augmentation, the fidelity metrics should
be interpreted as degradation relative to the full-information posterior, not as
exact posterior error under a separately constructed missing-data reference
posterior.

\subsection{Calibration Metrics}
\label{sec:eval-calibration}

Calibration is assessed with two complementary diagnostics. \textbf{Simulation-based
calibration (SBC)} examines the ranks of true parameters within posterior samples;
well-calibrated posteriors yield approximately uniform rank histograms.
\textbf{TARP} evaluates posterior calibration through expected coverage probability
curves; GAPSBI stores TARP plots and arrays along with aggregate deviations from
the identity coverage line (such as mean absolute error and integrated absolute
error).

\subsection{Posterior Fidelity Metrics}
\label{sec:eval-fidelity}

Posterior fidelity is measured by how closely a method posterior matches the
task-specific high-quality reference posterior for the corresponding reference
observation.
Let $P$ denote the method posterior and
$Q$ the reference posterior for a given observation, with sample sets
$\{\mathbf{a}_i\}_{i=1}^{n} \sim P$ and $\{\mathbf{b}_j\}_{j=1}^{m} \sim Q$.

\textbf{Classifier two-sample test (C2ST).} Following the SBIBM convention, an
MLP classifier $h$ is trained to distinguish method posterior samples from
reference posterior samples. Both sample sets are z-scored using the reference
samples. For computational tractability in repeated seed and budget sweeps, the
benchmark uses two hidden layers of width $5d_\theta$ and 3-fold
cross-validation, with 2,000 matched samples per posterior. The test statistic
is the mean cross-validation accuracy,
\begin{equation}
  \mathrm{C2ST} =
  \frac{1}{|\mathcal{T}|} \sum_{(\mathbf{z}, y) \in \mathcal{T}}
  \mathbb{1}\!\left[ h(\mathbf{z}) = y \right].
  \label{eq:c2st}
\end{equation}
Accuracy near $0.5$ indicates indistinguishable posteriors, while higher accuracy
indicates greater divergence. C2ST is computed separately for each fixed
reference observation using matched sample counts.

\subsection{Posterior Shift Metrics}
\label{sec:eval-shift}

To characterize \emph{how} a method posterior differs from its high-quality
reference, GAPSBI decomposes the shift into location and uncertainty-volume
components using moment-based metrics, computed per observation. Let
$\boldsymbol{\mu}_{\mathrm{method}}, \Sigma_{\mathrm{method}}$ and
$\boldsymbol{\mu}_{\mathrm{ref}}, \Sigma_{\mathrm{ref}}$ denote the empirical
mean and covariance of the method and reference posterior samples,
respectively.

The \textbf{Euclidean posterior mean shift} is the distance between posterior
means,
\begin{equation}
  D_E = \lVert \boldsymbol{\mu}_{\mathrm{method}} - \boldsymbol{\mu}_{\mathrm{ref}} \rVert_2 .
  \label{eq:euclidean}
\end{equation}
Values near $0$ indicate little location shift. The \textbf{covariance trace
ratio} compares average marginal variances,
\begin{equation}
  R_{\mathrm{tr}} =
  \frac{\operatorname{tr}(\Sigma_{\mathrm{method}})}{\operatorname{tr}(\Sigma_{\mathrm{ref}})},
  \label{eq:traceratio}
\end{equation}
with values above $1$ indicating larger average marginal variance. The
reference metrics are computed from stored posterior samples without retraining.
Together these separate posterior bias (Equation~\eqref{eq:euclidean}) from
posterior broadening (Equation~\eqref{eq:traceratio}).

\subsection{Variability Baselines}
\label{sec:eval-variability}

Because NPE training is stochastic, some posterior shift is expected even between
runs trained on identical complete data with different seeds. GAPSBI therefore
computes full-data variability baselines: it compares full-data runs against the
same high-quality reference posteriors using C2ST (Equation~\eqref{eq:c2st}) and
the same per-observation moment metrics
(Equations~\eqref{eq:euclidean}--\eqref{eq:traceratio}).
These baselines establish the
intrinsic seed-to-seed variability of full-data posteriors, against which
missing-data posterior shifts can be interpreted, so that differences attributed
to missingness are distinguished from differences attributable to training
stochasticity.

\section{Results}
\label{sec:results}

\textit{Placeholder.} This section will report empirical results once finalized.
Planned structure:
\begin{itemize}
  \item Calibration (SBC, TARP) across tasks, missingness mechanisms, and
  fractions.
  \item Posterior fidelity (C2ST) of each method relative to the fixed
  high-quality reference posteriors.
  \item Posterior shift decomposition (mean shift vs.\ covariance changes).
  \item Comparison against full-data seed variability baselines.
  \item Effect of simulation budget (high vs.\ mid vs.\ low).
\end{itemize}

\section{Discussion}
\label{sec:discussion}

\textit{Placeholder.} This section will interpret the results, including how
missingness mechanism and fraction interact with method behavior and calibration,
and what the posterior shift decomposition reveals about information loss versus
posterior displacement.

An important design choice in GAPSBI is to fix the inference backbone while
varying the missing-data representation. This makes conclusions about
representations stronger: if two methods differ mainly in whether they impute,
append the mask, embed observed coordinates, or learn a completion model, then
differences in calibration or posterior shift are more directly attributable to
that representation. Future versions of GAPSBI can naturally incorporate
architecturally distinct SBI methods, including transformer-based or
foundation-model-based posterior estimators, but this paper intentionally limits
architectural variability to isolate the question of representation under a
common NPE backbone.

\section{Limitations}
\label{sec:limitations}

\textit{Placeholder.} Anticipated limitations to be discussed include the
synthetic and relatively low-dimensional nature of the current tasks, the focus on
a controlled subset of tasks and methods, and the scope of the
evaluation protocol. This section will be completed alongside the results.

The benchmark currently fixes the NPE architecture for the main representation
comparisons. This improves experimental control, but it does not answer how
fundamentally different SBI architectures compare under missing data. Evaluating
methods such as Simformer within the same dataset and evaluation protocol is a
natural direction for future work.

\section{Conclusion}
\label{sec:conclusion}

\textit{Placeholder.} This section will summarize the benchmark and its intended
role for studying simulation-based inference under missing data, and outline
directions for extension.

\bibliography{iclr2026_conference}
\bibliographystyle{iclr2026_conference}


\end{document}
